\chapter{점수와 채점: 무엇으로 고르고 무엇으로 평가하나}

5장의 $p$는 ``누가 더 위험한가''를 알려 준다. 이 장은 그 $p$로 \textbf{어떻게 고르는지}, 고른 결과를 \textbf{어떻게 채점하는지}를 다룬다. 고르는 도구와 채점하는 도구를 섞으면 결과가 무의미해지므로, 둘을 분명히 나눈다.

\section{기대수익 점수: $p$를 돈으로 바꾸기}
\begin{formulabox}[기대수익 점수]
\[ s=(1-p)\times r - p\times \mathrm{LGD} - tb \]
$r$: 약정 금리, LGD: 부도 시 손실률(그 시점까지 확정된 부도 대출에서 ``1 $-$ 받은 총액/대출액''의 평균, 2013년 모델 0.62, 2016년 0.52), $tb$: 같은 달 같은 만기 국채 수익률.
\end{formulabox}
\begin{itemize}[leftmargin=*]
\item $(1-p)\times r$: 안 망하면 받는 이자의 기댓값.
\item $p\times$LGD: 망하면 잃는 원금의 기댓값.
\item $tb$: 이 대출 대신 국채를 샀다면 위험 없이 받았을 수익(기회비용).
\end{itemize}

\begin{workout}[금리와 위험을 함께 봐야 하는 이유]
\Use{기대수익 점수(이 절)}
\Step{1} 예시: LGD 0.6, 국채 1.5\%. 대출 가($p$ 5\%, 금리 8\%), 대출 나($p$ 25\%, 금리 22\%).
\Step{2} 가: $0.95\times8\%-0.05\times0.6-1.5\%=7.6-3.0-1.5=+3.1\%$p.
\Step{3} 나: $0.75\times22\%-0.25\times0.6-1.5\%=16.5-15.0-1.5=0$.
\Step{4} 나의 실제 위험이 15\%라면: $0.85\times22\%-0.15\times0.6-1.5\%=18.7-9.0-1.5=+8.2\%$p.
\Ans{부도확률만 보면 항상 가를 사지만, 위험에 비해 금리가 후한 대출은 나 같은 고금리 대출일 수 있다.}
\Trap{점수를 ``정확한 예상 수익''으로 읽는다. $p$는 생애 기준이고 금리는 연 단위라, 이 식은 순위를 매기는 1년 단위 근사다.}
\end{workout}

점수는 두 곳에 쓴다. (1) 승인 규칙: $s>0$, 즉 국채보다 더 벌 것 같으면 승인. (2) 순위: 점수 상위 10\%, 20\%, $\dots$를 승인해 본다. 같은 달 같은 만기 대출은 국채가 같으므로 순위를 정하는 것은 사실상 ``금리와 위험의 균형''이다.

\section{고르는 도구와 채점하는 도구}
\begin{table}[H]\centering\small
\begin{tabularx}{\linewidth}{@{}p{24mm}LL@{}}\toprule
 & 기대수익 점수 & IRR(현금흐름 수익률)\\\midrule
언제 & 승인하는 날 & 대출이 끝난 뒤\\
무엇으로 & 모델의 $p$, 약정 금리, 평균 손실률, 국채 & 실제로 받은 돈과 받은 시점, 수수료 1\%\\
정체 & 예상 & 실제 결과\\
역할 & 줄 세워 고른다 & 고른 묶음이 실제로 얼마 벌었는지 채점\\\bottomrule
\end{tabularx}
\caption{두 도구의 역할(2장의 IRR 참고).}
\end{table}
IRR로 고를 수는 없다(승인하는 날에는 미래를 모른다). 점수로 채점할 수도 없다(자기 예상으로 자기를 채점하는 셈이다). \textit{비유:} 학원 예상 점수로 지원 대학을 고르고, 실제 수능 성적으로 그 예상이 맞았는지 채점한다.

\section{왜 ``아무거나 고를 때''와 비교하나}
\textbf{질문.} 우리가 고른 대출의 수익률이 7\%였다. 잘한 걸까?
\begin{itemize}[leftmargin=*]
\item 그해 아무거나 샀어도 6\%였다면 우리 실력은 1\%p다.
\item 그해 아무거나 샀어도 8\%였다면 우리는 고르다가 1\%p를 잃었다.
\end{itemize}
수익률 하나로는 구분이 안 된다. 그래서 같은 달, 같은 승인비율로 \textbf{무작위로} 고른 결과를 빼서 ``고르는 실력''만 남긴다. 펀드매니저를 시장 지수 대비로 평가하는 것과 같은 원리다.

\begin{formulabox}[선별 가치]
\[ \mathrm{SV}=\frac19\sum_{q=10\%}^{90\%}\Big[\overline{(\mathrm{IRR}-tb)}_{\text{점수 상위 }q}-\overline{(\mathrm{IRR}-tb)}_{\text{무작위 }q}\Big] \]
승인 대상은 발행월마다 그 달 대출 중 상위 $q$로 자르고, 발행월 묶음 IRR을 계산해 평균한다. 승인비율 하나를 골라 유리하게 보일 수 없도록 9개 비율을 평균한다. (25단계까지는 발행연도마다 잘랐다. 발행월로 바꿔도 1.03 \ensuremath{\to} 1.07\%p로 거의 같다.)
\end{formulabox}

이 기준이 좋은 이유는 세 가지다. (1) 시장 효과가 양쪽에 똑같이 들어가 지워진다. (2) 승인비율을 맞추므로 ``얼마나 샀나''가 아니라 ``무엇을 샀나''만 비교한다. (3) 0이 뜻을 갖는다. 무작위는 정보가 없는 투자자이고, 실제로 정답을 섞은 가짜 모델은 0 근처였다.

\begin{table}[H]\centering\small
\begin{tabularx}{\linewidth}{@{}p{34mm}LL@{}}\toprule
비교 대상 & 답하는 질문 & 주 기준이 아닌/인 이유\\\midrule
국채 & LC 대출에 투자할 가치가 있나? & 2013--16년은 전부 사도 국채를 이겼다. 모델 없이도 양수라 고르는 실력을 못 보여 준다\\
전부 승인 & 모델을 쓰는 게 나은가? & 승인비율이 달라 ``무엇''과 ``얼마나''가 섞인다\\
\textbf{같은 비율 무작위} & \textbf{고르는 방식이 가치를 더하나?} & \textbf{주 기준. 시장과 승인비율을 모두 맞춘다}\\\bottomrule
\end{tabularx}
\end{table}

\section{Sharpe 이야기: 세 가지 정의}
교수님은 ``결과는 Sharpe 극대화''를 말씀하셨고, 조교는 ``0.2--0.4면 옳게 한 것''이라고 했다. 그런데 Sharpe는 무엇을 한 단위로 보느냐에 따라 전혀 다른 숫자가 된다.

\begin{table}[H]\centering\small
\begin{tabularx}{\linewidth}{@{}p{40mm}LLp{20mm}@{}}\toprule
정의 & 계산 & 특징 & 우리 값\\\midrule
건별 Sharpe(거절 = 0) & 모든 대출에 대해, 승인이면 그 대출의 초과수익, 거절이면 0. 평균 \ensuremath{\div} 표준편차 & 승인비율만 줄여도 변한다(아래) & 0.03--0.15\\
\textbf{건별 Sharpe(승인만)} & 승인한 대출들의 초과수익 평균 \ensuremath{\div} 표준편차 & \textbf{조교의 0.2--0.4와 크기가 맞는다} & 0.10--0.34\\
발행월 포트폴리오 Sharpe & 매달 (승인 대출 + 거절 자금 국채) 포트폴리오의 초과수익을 48개월에 걸쳐 평균 \ensuremath{\div} 표준편차 & 승인비율이 지워진다. ``매달 얼마나 고르게 벌었나''를 잰다 & 약 2--3.5\\\bottomrule
\end{tabularx}
\caption{Sharpe 정의 세 가지(32단계). 건별 정의의 대출 수익률은 받은 돈을 국채로 재투자한 만기 수익률이다(대출 한 건의 IRR은 부도 시 $-77\%$처럼 극단값이 커서 쓰지 않는다).}
\end{table}

\begin{workout}[거절을 0으로 세면 생기는 일]
\Use{건별 Sharpe(거절 = 0)}
\Step{1} 예시: 승인한 대출의 초과수익 평균 4\%, 표준편차 20\%. 승인비율 $a$만 바꾼다.
\Step{2} 승인 100\%: $0.04/0.20=0.200$.
\Step{3} 승인 50\%: $\sqrt{0.5}\times0.04/\sqrt{0.04+0.5\times0.0016}=0.140$.
\Step{4} 승인 10\%: 0.062.
\Ans{대출의 질은 같은데 숫자만 1/3이 된다. 이 정의는 ``얼마나 승인했나''에 끌려간다.}
\end{workout}

포트폴리오 Sharpe는 이 문제가 없다(팀원 제안서의 지적, 30단계). 매달 초과수익이 (승인 금액 비중 \ensuremath{\times} 승인 묶음 초과수익)이므로 비중이 평균과 표준편차에서 함께 지워진다. 다만 규모를 무시하고, 위험을 2013--16년(경기침체 없음) 48개월의 출렁임으로만 재며, 그 크기(약 3)는 주식 Sharpe(0.5 안팎)와 비교할 수 없는 \textbf{순위 비교용} 숫자다.

\begin{warnbox}[자르는 단위가 중요하다]
포트폴리오 Sharpe는 매달 승인 금액 비중이 일정하다는 전제에서 승인비율이 지워진다. 승인 대상을 발행연도 단위로 자르면 고금리 대출이 많은 달에는 많이, 적은 달에는 적게 승인되어 비중이 들쭉날쭉해지고, 그 출렁임이 Sharpe를 깎는다. 실제로 기대수익 순의 Sharpe는 연도 단위 2.46, 발행월 단위 2.91이었다(외부 검증, 37단계). 연도 단위는 12월 대출까지 본 뒤 1월 승인을 정하는 셈이라 실전에서도 쓸 수 없다.
\end{warnbox}

\section{그래서 두 축으로 본다}
\begin{keybox}[수익 축과 안정성 축]
\begin{itemize}[leftmargin=*]
\item \textbf{수익:} 선별 가치(IRR 기준), 총 NPV(국채로 할인한 순현재가치, 달러 단위).
\item \textbf{안정성:} 발행월 포트폴리오 Sharpe, 승인 대출 건별 Sharpe, 하위 5\% 평균(꼬리 손실).
\end{itemize}
Sortino(아래쪽 출렁임만 위험으로 보는 Sharpe)도 계산했지만, 9개 승인비율 모두에서 전략 순위가 건별 Sharpe와 같아 새 정보가 없었다(33단계). 포트폴리오 단위 Sortino는 48개월 모두 국채보다 높아 분모가 0이 되어 계산할 수 없다.
\Intuit{더 버는 규칙과 더 고르게 버는 규칙은 다를 수 있다. 하나의 숫자로 우열을 정하지 말고 두 축으로 본다.}
\end{keybox}

\section{성과의 출처를 나누는 두 도구}
\paragraph{Brinson 분해.} 펀드 평가 도구다. 초과수익 = 배분(좋은 등급을 더 많이 샀나) + 선별(같은 등급 안에서 더 나은 대출을 골랐나) + 상호작용. 예를 들어 기준이 A 60\%(4\%)·\allowbreak{}B 40\%(7\%)로 5.2\%이고 모델이 A 70\%(5\%)·\allowbreak{}B 30\%(8\%)로 5.9\%라면, 0.7\%p = 배분 $-0.3$ + 선별 $+1.0$ + 상호작용 0이다.

\paragraph{등급 중립 롱숏.} 주식 팩터 투자에서 신호의 크기를 재는 방법이다. 매달, 등급마다 신호 상위 30\%를 산다고(롱) 하위 30\%를 판다고(숏) 가정하고 두 묶음의 IRR 차이를 기록한다. 등급마다 롱과 숏을 같은 비율로 잡으므로 등급 배분과 시장 효과가 정확히 지워지고 \textbf{등급 안 선별}만 남는다. LC 대출은 공매도가 불가능하므로 실제 전략이 아니라 \textbf{측정 도구}다. 실제로 가능한 형태는 ``상위 30\% $-$ 그 등급 전체''(롱온리)다.
